\cotefile{101}
\scanpage{16}
algébriques, \unsure{en} mettant sur $H^0(Y,G)\simeq\deleted{\unsure{H}}\prod_i\unsure{\mathcal O|\mathscr O}_{x_i}$ la topologie
\deleted{produit des topologies} naturelle de l'anneau \unsure{$\frac12$ local|semi-local|local},
[topologie \unsure{$J$-adique|fonctionnelle} en \unsure{fait|produit}], et sur $H^1(U;F)$ une
topologie \deleted{\unsure{extérieure|quotient}} \unsure{s.:|et} \deleted{\illegible{1}}\inserted{la} discrète \deleted{\unsure{sur}} sur $I$ pour la
topologie quotient sur $\prod_i\mathcal O_{x_i}$ [N.B., $A$ est un
\unsure{sous-groupe|sous-espace} \uline{complet}\inserted{\unsure{donc} \uline{fermé}}, donc le quotient $\prod_i\mathcal O_{x_i}/A$
est séparé].

Question : \unsure{classifier|dualité} \mbox{-{}-{}-{}-}

Autres questions : $\operatorname{Ext}^i(K,A)=0$ pour tt $i$ ($A$ anneau local
\unsure{noeth|noeth.} \unsure{régulier|intègre} complet, de corps $K$)
% archive sheet number: 15

\scanpage{17}
Notes antiques \hfill 1958

(Caractérisation des
\unsure{Coh. M.|Coh. Mc.|Coh. Mac.}, des
\unsure{Gorensteins|Gorenstein} etc)
% archive sheet number: 16

\scanpage{18}
\S\,1 : \uline{Les foncteurs} $H^i_{\unsure{\mathscr C|\tau|\mathfrak C}}$

\noindent\makebox[0pt][r]{\smash{\raisebox{-12\baselineskip}{\(\left\{\rule{0pt}{24\baselineskip}\right.\)}}\hspace{0.6em}}%
1. Sur les foncteurs $A$-linéaires (de modules et modules topologiques)

2. Définition des $H^i_{\mathscr C}(M)$, premières propriétés.

3. Fonctorialité pour les $H^i_{\mathscr C}(M)$ \inserted{(y inclure \unsure{cup-produit|cup-produits})}

4. Cas d'un anneau local $A$ : \inserted{premières propriétés spéciales; \uline{caractérisation des} $H^{\unsure{i}}$}

\deleted{\unsure{4}}5. \inserted{relations avec la codim. coh.} Rappels sur les Modules de Cohen-Macaulay

6. Résultats auxiliaires sur certains $\operatorname{Ext}_A^i(N,M)$ \margin{[$M$ \unsure{Ch. Mac|Coh. Mac.} de supp. $\operatorname{Spec}(A)$}

\deleted{\unsure{6}}7. Cas particuliers

\S\,2 \uline{Foncteurs et modules dualisants} \inserted{(théorie locale)} \uline{Le th. de dualité locale}.

1. Définition et premières propriétés des foncteurs et modules dualisants.

2. Extension du foncteur dualisant

3. Th. d'existence et unicité; exemples.

4. \deleted{\unsure{Le|Les|La} \unsure{théorème|théorèmes|théorie} de dualité locale} Structure des $H^i(M)$.

5. Module fondamental et \uline{th. de dualité locale}

6. Anneaux de Gorenstein.


\S\,3 \uline{Les Foncteurs} $H_i$, \uline{et les complexes dualisants}
% archive sheet number: 17

\scanpage{19}
\margin{\unsure{1}\quad\unsure{?}}
\deleted{\unsure{Revenons|Remarque} : \unsure{à la} question de dualité. Cela \unsure{veut} \unsure{dire}
la validité du (iii) du lemme \unsure{3|2} signifie
que \unsure{la dimension} $\dim A\leq n$. \unsure{et que} \illegible{2}
\unsure{trouve que|donne}}

\uline{Théorème 4}

\uline{Définition} Soit $A$ anneau \deleted{local} noethérien de dim $n$.
Un module $\Omega$ sur $A$ est dit un \textquotedblleft module fondamental\textquotedblright{}
\deleted{\unsure{coefficient de dualité}} si il satisfait les deux conditions

% slanted
\margin{Introduire $\Omega$ par une proposition, disant que (i) et (ii) équivalent à certaines propriétés du foncteur $M\to\operatorname{Hom}(M,\Omega)$.}

(i) $\Omega$ est de type fini

(ii) $\deleted{\unsure{H^n(\Omega)}} D(\Omega)\simeq H^n(A)$

Un tel $\Omega$ est unique à isom. près \deleted{\unsure{au|si}} \inserted{\uline{si}}
$A$ est local (ou semi-local), \unsure{puis que|puisque|mais que} $\deleted{\unsure{D(A)}=}\inserted{\widehat\Omega=D(H^n(A))}$
est alors \unsure{connu|unique}. On le note $\Omega_A$. D'autre part si $A$ est \inserted{local} complet, un tel $\Omega_A$ existe
\unsure{tjrs} (on prend $D(H^n(A))$). Il existe aussi
si $A$ régulier \inserted{et équidimensionnel} : prendre $\Omega=A$ !

\deleted{D'ailleurs, D'après le \unsure{th'|th.} 3,
l'existence de $\Omega$ implique que \unsure{est} \unsure{lgs de un seul}
\illegible{1} \unsure{annulateur|accouplement} de \unsure{ce} $H^n(A)=$}

\margin{Si $A$ est artinien, alors \uline{module fond} $=$ \uline{module dualisant}}

\deleted{\uline{Proposition} Soit $A$ un anneau local}

\uline{Théorème 4} Soit $A$ \inserted{\unsure{locl|local}} noethérien\inserted{de dim $n$}, muni d'un
module \deleted{dualisant} \inserted{\uline{fondamental}} $\Omega$. Le \unsure{.} \unsure{un} \deleted{\unsure{isom.}}\inserted{hom}
\[
H^n(\Omega)\simeq H^n(A)\otimes\Omega\xrightarrow{\unsure{g|\eta}} \mathrm I
\]
[déduit d'un isom $H^n(A)\xrightarrow{\sim}D(\Omega)$], définit
des \unsure{accouplements|accouplement} \unsure{donnant} $\unsure{\alpha_i}$, qui est un
isom en dim $n$. Pour que ce soit un
% archive sheet number: 18

\scanpage{20}
\deleted{\uline{Corollaire 1 du th. 3} \inserted{Si $A$ est de dim $n$,} Le foncteur $M\unsure{-|\mapsto}H^n(M)$ est
exact à droite, et \deleted{\unsure{est}}\inserted{\unsure{°}} par suite}
\[
\deleted{\boxed{H^n(M)\xleftarrow{\sim}H^n(A)\otimes M}}
\]

\deleted{Corollaire 3 Si $D(\Omega)\simeq H^n(A)$, alors $H^n(\Omega)\simeq D(\operatorname{Hom}(\Omega,\Omega))$}

\inserted{\uline{Corollaire 2} \deleted{\unsure{Soit $M$}} \unsure{de} dim $q$, on a $\operatorname{Ext}^q(k,M)\ne0$. Si $i>q$ est tel que
$\operatorname{Ext}^i(k,M)\ne0$, alors $\operatorname{Ext}^j(k,M)=0$ pour $j>i$.}

\deleted{Soit \unsure{alors|donc} $\Omega$ donné, et $H^n(\Omega)\simeq\deleted{\unsure{D(\Omega)}}H^n(A)\otimes\Omega\xrightarrow{\varphi}\mathrm I$}

\deleted{\unsure{ce qui revient}}

\deleted{Il revient au \unsure{$=$|même} de se donner $\Omega\xrightarrow{\varphi_2}D(H^n(A))$
ou $H^n(A)\xrightarrow{\varphi_1}D(\Omega)$. On a alors des
accouplements}
\[
\deleted{H^p(M)\times\operatorname{Ext}^{n-p}(M,\Omega)\longrightarrow H^n(\Omega)\longrightarrow\mathrm I}
\]
\deleted{d'où}
\[
\deleted{H^p(M)\longrightarrow D\bigl(\operatorname{Ext}^{n-p}(M,\Omega)\bigr)}
\]

\deleted{Explicitons-les en dims extrêmes.}
\[
\deleted{H^n(M)\simeq\deleted{\unsure{M\otimes}}H^n(A)\otimes M\longrightarrow D(\operatorname{Hom}(M,\Omega))\simeq D(\Omega)\otimes M}
\]
\deleted{\deleted{est \unsure{précisément associé à} $\unsure{\varphi_1}$}
s'identifie à $\varphi_1\otimes\operatorname{id}_M$ ; pour que
ce soit un isom pour tt $M$, il f. et s.
que $\varphi_1$ soit un isom. Considérons le
cas \unsure{part.} \deleted{$\unsure{\varphi_1\otimes(M)}$} $M=\Omega$, on trouve}

% slanted
\margin{plus bas ?}

\[
\deleted{\deleted{H^n(\Omega)\unsure{\simeq} H^n(A)}}
\]
\[
\deleted{H^n(\Omega)\times\operatorname{Hom}(\Omega,\Omega)\longrightarrow H^n(\Omega)\longrightarrow\mathrm I}
\]
\deleted{d'où}
\[
\deleted{\widehat{\operatorname{Hom}(\Omega,\Omega)}\xrightarrow{\sim}D(H^n(\Omega))}
\]
\deleted{transformant l'identité en $\varphi$. \deleted{on} \uline{\unsure{donc|dira}}
\unsure{on a} équivalence des conditions suivantes :}

\deleted{(i) $H^n(\Omega)$ est dualisant et \unsure{que} \illegible{1}}

\deleted{(ii) $\operatorname{Hom}(\Omega,\Omega)\simeq A$ \qquad (iii)}
% archive sheet number: 19

\scanpage{21}
% vertical rule at left of the continuation
\noindent\llap{\smash{\raisebox{-0.7\baselineskip}{\rule{0.4pt}{1.6\baselineskip}}}\quad}
isom pour les dim. $p\geq k$ ($k$ entier donné $\leq n$)
il faut et il suffit qu'on ait
\[
H^i(A)=0\qquad\text{pour}\quad k\leq i<n
\]

% vertical rule at left of the corollary
\noindent\llap{\smash{\raisebox{-2.1\baselineskip}{\rule{0.4pt}{2.2\baselineskip}}}\quad}
\uline{Corollaire 1} Pour que ce soit un isom. pour
\deleted{\uline{tout}} $p$, il f. et s. que \inserted{les} $A_{\unsure{m|\mathfrak m}}$ \unsure{soient|soit} \deleted{des}
\deleted{\unsure{modules}} \inserted{anneaux} de Coh. Mac. de dim $n$.

\deleted{\uline{\unsure{Corollaire}}}

\deleted{\uline{Proposition} \unsure{9|2}}\inserted{\uline{Corollaire 2}} Soit $A$ un anneau \inserted{local} de Coh.
Mac. de dim $n$, $\Omega$ un module \deleted{dualisant}\inserted{fond.}
pour $A$. Alors on a
\[
\left\{\begin{aligned}
&\operatorname{Ext}^i(k,\Omega)=0\qquad \text{si}\quad i\ne n\\
&\operatorname{Ext}^n(\unsure{k|h},\Omega)\simeq k
\end{aligned}\right.
\]

\margin{\unsure{\S|$\{$}}
En effet, $\operatorname{Ext}^i(k,\Omega)$ est dual de $H^{n-i}(k)$, qui
est nul si $n-i\ne0$, et qui est $=\underline{k}$ si $n-i=0$

En particulier

\uline{Corollaire 3} $A,\Omega$ comme ci-dessus. Alors $\Omega$
est un module de Coh. Mac. \inserted{de dim $n$} sur $A$, on a
\[
\left\{\begin{aligned}
&H^i(\Omega_{\unsure{A|*}})=0\qquad \text{si}\quad i\ne n\\
&H^n(\Omega_A)\quad\text{est dualisant}
\end{aligned}\right.
\]
De plus, \deleted{$\Omega_A$ est un module fidèle}
\[
\left\{\begin{aligned}
&\operatorname{Ext}^i(\Omega,\Omega)=0\qquad\text{si }i\ne0\\
&\deleted{\operatorname{Ext}}\ \operatorname{Hom}(\Omega,\Omega)\simeq A
\end{aligned}\right.
\]
% archive sheet number: 20

\scanpage{22}
Que $\Omega$ soit Coh. Mac. de dim résulte des
formules précédentes et de la caractérisation des modules
de Coh. Mac. de dim $n$. Il s'ensuit
$H^i(\Omega)=0$ pour $i<n$, \deleted{\unsure{et}}\inserted{et pour $i>n$ d'après le th. 3.} D'après le th. 1,
le cor. 1 implique que $H^n(\Omega)$ est dualisant.

D'après le th. de dualité, on a
\[
H^i(\Omega)\xrightarrow{\sim}D(\operatorname{Ext}^{n-i}(\Omega,\Omega))
\]
-- qui montre que \deleted{$H^i(\Omega)=0$ si $i>n$.}
$\operatorname{Ext}^{n-i}(\Omega,\Omega)=0$ si $\deleted{\unsure{n-}}i\ne n$ i.e. $n-i\ne0$.
\deleted{\unsure{Cela montre aussi que}} Si $\deleted{\unsure{n-}}i=n$, on
trouve $D(\operatorname{Hom}(\Omega,\Omega))\simeq H^n(\Omega)$ qui est
dualisant i.e. qui est $\simeq D(A)$, d'où
$\operatorname{Hom}(\Omega,\Omega)\simeq A$. Ainsi \uline{tout endomorphisme
de $\Omega$ est donné par un scalaire}, déterminé de
\uline{façon unique}. \inserted{A fortiori\unsure{;|,}} $\Omega$ \uline{est \deleted{\unsure{donc}} un module fidèle}.

\uline{Corollaire 4} \inserted{$A,\Omega,n$ comme dessus.} $\Omega_{\unsure{A}}$ est de dim inj. $=n$, en
particulier on a
\[
\operatorname{Ext}^i(M,\Omega)=0\qquad\text{si}\quad i>n
\]
\uline{quel que soit} $M$

Il suffit de le vérifier pour $M$ de
type fini, et ça résulte alors du th. de dualité.

% margin note written sideways
\margin{\uline{Corollaire 5} Soit $A$ un anneau local de Coh. Mac. de dim $n$, $\Omega$ un module fondamental pour $A$, $M$ un $A$-module de type fini\unsure{.} Alors \deleted{\unsure{on a}} $\operatorname{cocodim}M\leq k\Longleftrightarrow\operatorname{Ext}_A^{\unsure{i}}(M,\Omega)=0$ pour $i>n-k$. Si $M$ est de dim Krull $k$, cela équivaut alors à : $M$ est Coh. Mac. et implique $\operatorname{Ext}_A^i(M,\Omega)=0$ pour $i\ne k$.}
% archive sheet number: 21

\scanpage{23}
$A$ anneau local de dim $n$, $\mathrm I$ un module dualisant
$\Omega$ un module fondamental pour $A$, avec un isom donné
\[
\boxed{\unsure{g|\eta|\mathfrak D}_1:H^n(A)\xrightarrow{\sim}D(\Omega)}
\]
correspondant à un accouplement
\[
\begin{aligned}[b]&\text{[où \deleted{\unsure{un}}}\end{aligned}\quad
\boxed{\begin{aligned}[b]
&H^n(A)\otimes\Omega=H^n(\Omega)\xrightarrow{g}\mathrm I\\
&g\in D(H^n(\Omega))
\end{aligned}}\quad\begin{aligned}[b]&\text{\deleted{\unsure{est}} ]}\end{aligned}
\]
et à un isom
\[
g_2:\widehat\Omega\xrightarrow{\sim}D(H^n(A))
\]

On a pour \uline{tout $M$ de type} fini
\[
\boxed{H^n(M)\xleftarrow[\sim]{\mathrm{can}}H^n(A)\otimes M}
\xrightarrow[\sim]{g_1\otimes\operatorname{id}_M}D(\Omega)\otimes M
\xrightarrow[\sim]{\mathrm{can}}D(\operatorname{Hom}(M,\Omega))
\]
d'où
\[
\begin{aligned}[b]&\text{et par là}\end{aligned}\quad
\boxed{\begin{aligned}[b]
g_{1,M}:{}&H^n(M)\xrightarrow{\sim}D(\operatorname{Hom}(M,\Omega))\\
g_{2,M}:{}&\operatorname{Hom}(M,\Omega)^{\wedge}\xrightarrow{\sim}D(H^n(M))
\end{aligned}}
\]
qui pour $M\simeq A$ se réduisent à $g_1,g_2$. Pour $M=\Omega$, on trouve
\[
\begin{aligned}
g_{1,\Omega}:{}&H^n(\Omega)\xrightarrow{\sim}D(\operatorname{Hom}(\Omega,\Omega))\\
g_{2,\Omega}:{}&\operatorname{Hom}(\Omega,\Omega)^{\wedge}\xrightarrow{\sim}D(H^n(\Omega))
\end{aligned}
\]
Enfin, on constate que l'élément identique de
$\operatorname{Hom}(\Omega,\Omega)^{\wedge}$ correspond à l'élément $g$ de $D(H^n(\Omega))$.
Par suite, le transposé de $g$ est l'hom
\[
{}^tg:\widehat A\longrightarrow\operatorname{Hom}(\Omega,\Omega)^{\wedge}
\]
qui transforme $1_{\widehat A}$ en $\operatorname{id}_{\widehat\Omega}$, donc $\lambda$ en l'homothétie $\lambda$.

\uline{Proposition} 9 L'annulateur de $\Omega$ est l'intersection $\unsure{\mathfrak N|N}$
des idéaux primaires associés aux idéaux pr. de
corang $n$ dans $A$, et $\Omega$ est aussi un
module fondamental pour $A'=A/\mathfrak N$. Les idéaux
pr. associés à $\Omega$ sont les idéaux pr. de corang $n$ de $A$. \unsure{$p$|$\mathfrak p$}
\margin{La longueur $\unsure{\mathrm{lg}_p|\mathrm l_p|\ell_{\mathfrak p}}(\Omega)=\mathrm{lg}_p(A)$ pour un tel $p$}

L'annulateur de $\Omega$ est celui de $D(\Omega)=H^n(A)$,
donc \deleted{\unsure{notre|toute}} la première assertion suit du th. 3; la
% archive sheet number: 22

\scanpage{24}
dernière suit du cor. 1 du dit. On a
\[
0\unsure{-|\to}\unsure{\mathfrak N|N}\to A\to A'\to0
\]
d'où, comme $\mathfrak N$ est de dim $<n$ donc $H^n(\mathfrak N)=0$,
$H^n(A)\simeq H^n(A')$ donc $\Omega$ est aussi \unsure{fond|fond.} pour $A'$.

[N.B. $\mathfrak N$ est le plus grand idéal par lequel on
peut diviser \unsure{sans} que $\Omega$ cesse de \unsure{être|rester} \unsure{fond}]

\uline{Corollaire 1} L'image de $H^n(\Omega)\xrightarrow{\unsure{g|\eta|\mathfrak D}}\mathrm I$ est l'annulateur
de $\mathfrak N$, \inserted{\unsure{;|i.e.}\deleted{\unsure{et}} $\operatorname{Coker}g\simeq D(\mathfrak N)$.} \deleted{Donc $g$ est surjectif si et seulement
si \unsure{tous les} \unsure{les idéaux}} \inserted{Les conditions suivantes sont équiv.}

(i) $g$ est surjectif

(ii) $\mathfrak N=0$, i.e les idéaux premiers associés à $0$ de $A$
sont tous de \deleted{dim} corang $n$

(iii) $\Omega$ est un $A$-module fidèle.

En effet, la première assertion résulte du
fait que $\mathfrak N$ est le noyau de $A\to\operatorname{Hom}(\Omega,\Omega)$,
qui est transposé de $g$. D'où les équivalences.

\uline{Corollaire 2} $\operatorname{Ker}g\simeq D\bigl(\operatorname{Hom}(\Omega,\Omega)/A\,\operatorname{id}_{\Omega}\bigr)$.

Les conditions suivantes sont équivalentes

(i) $g$ est injectif \qquad (i bis) $H^n(\Omega)$ est comonogène

(ii) Tout end. de $\Omega$ est une homothétie

(ii bis) $\operatorname{Hom}(\Omega,\Omega)$ est monogène

La première formule est de dualité.
$(\mathrm i)\Rightarrow(\mathrm{i\,bis})$, \deleted{\unsure{$(\mathrm{ii})$}} $(\mathrm{i\,bis})\Longleftrightarrow(\unsure{\mathrm{ii}'|\mathrm{ii}}\,\mathrm{bis})$ et
% archive sheet number: 23

\scanpage{25}
(ii) $\longrightarrow$ (ii bis) triviaux. Enfin \deleted{\unsure{(i bis)}}\inserted{(ii bis)} $\Rightarrow$ (ii), car
si $M$ est un module de type fini sur $A$ tel que
$\operatorname{Hom}(M,M)$ soit monogène, il résulte de Nakayama \unsure{que}
$\operatorname{id}_M$ en est un générateur.

\uline{Cor. 3} Conditions équivalentes

(i) $\unsure{g|\eta|\mathfrak D}$ est bijectif

(ii) $H^n(\Omega)$ est dualisant

(iii) $A\to\operatorname{Hom}(\Omega,\Omega)$ est un isom

(iv) $\operatorname{Hom}(\Omega,\Omega)\simeq A$

\deleted{\unsure{et}} Alors \deleted{\unsure{aussi}} les idéaux pr. associés à $A$
sont tous de \deleted{\unsure{dim|codim}} corang $n$.

\uline{Cor. 4} Les cond. du cor. 3 sont vérifiées
sous les conditions suivantes : a) $A$ est
Cohen-Macaulay \qquad (b) $A$ est normal.

\deleted{\unsure{(a)}} (a) Résulte du th. 4, cor. 3.

% slanted
\margin{à expliciter en \unsure{lemme}}
Quant à (b), comme les conditions du cor. 1
sont vérifiées ($A$ est intègre) il suffit
de prouver que $\operatorname{Hom}(\Omega,\Omega)\unsure{>|\supset} A$ est en
fait $=A$. Or $\Omega$ est sans torsion
[\deleted{ses idéaux pr. associés} $0$ est le seul idéal pr. associé,
car $A$ intègre] donc $\operatorname{Hom}(\Omega,\Omega)\subset K$ corps
% archive sheet number: 24

\scanpage{26}
des fractions de $A$. Comme c'est un module de
type fini sur $A$, il est dans la clôture
intégrale de $A$, O.K.

\uline{Question} Les conditions du cor. 2 sont elles
\uline{tjrs} vérifiées ? On est ramené évid. au cas
où les idéaux pr. associés à $0$ dans $A$
sont de \deleted{dim} corang $n$ : est-il vrai alors
que $H^n(\Omega)$ est dualisant ? \uuline{Non} déjà
dans le cas \deleted{d'une surface \unsure{non normale} ;} \inserted{\unsure{voir} \deleted{\unsure{l'ex.}} un anneau}
local \inserted{intègre} $A$ [de dim 2], tel que $A_p$ soit
régulier pour $p$ \deleted{\unsure{max}}\inserted{premier \unsure{minimal|minimum}} [$A$ non normal]. En
effet, \deleted{\unsure{on vérifie que}} on prouve alors que
\[
\operatorname{Hom}_A(\Omega,\Omega)=\bigcap_{\substack{p\ \text{premier}\\\text{de rang }1}}A_p
=\deleted{\text{\unsure{normalisé}}}\inserted{\text{clôture intégrale}}\text{ de }A .
\]
[$\Omega$ est en effet défini par des conditions \unsure{valuatives}
\mbox{-{}-{}-} ].

\uline{Proposition 10} Soit $x$ un élément de $A$ qui
n'y est pas diviseur de $0$, et tel que
$xH^{n-1}(A)=H^{n-1}(A)$. Alors $\Omega_A/x\Omega_A$ est
un module fondamental pour $\deleted{\unsure{\Omega_A|\Omega}}A/xA$
% archive sheet number: 25

\scanpage{27}
En effet, de la suite exacte
\[
0\unsure{-|\to} A\xrightarrow{x}A\to A/xA\to0
\]
on tire
\[
H^{n-1}(A)\xrightarrow{x}H^{n-1}(A)\to H^{n-1}(A/xA)\to H^n(A)\xrightarrow{x}H^n(A)
\]
et compte tenu de l'hypothèse
\[
H^{n-1}(A/xA)\simeq\deleted{\operatorname{Ker}\text{ de }x\text{ dans}}\inserted{\text{ann. de }x}\text{ dans }H^n(A)
\]
\deleted{\unsure{a}}
Comme $H^n(A)\simeq D(\Omega)$, on trouve bien
\[
H^{n-1}(A/xA)\simeq D(\Omega/x\Omega),\qquad\text{cqfd.}
\]

\uline{Corollaire} Supposons $A$ Coh. Mac. et soit
$x_1\mathord{-}\,\mathord{-}x_p$ une partie d'un syst. de
paramètres de $A$. Alors $\Omega_A/(x_1,\mathord{.}\,\mathord{-},x_p)\Omega_A$ est
un module fondamental pour $A/(x_1,\mathord{-},x_p)$.

Supposons maintenant que $A$ soit un
quotient d'un anneau local \deleted{\uline{\unsure{régulier}}}\inserted{Coh. Mac.} $A'$
\deleted{\unsure{par}}\inserted{de dim \unsure{$n'$|$N'$}} 
\[
A\simeq A'/\unsure{\mathfrak N|N}
\]
Supposons $A'$ muni d'un $\Omega_{A'}$. Si
$M$ est un $A$-module, c'est un $A'$-module,
et $D(M)$ peut-être interprété ainsi
[choisissons un module dualisant $I'$ \inserted{$=H^{\unsure{n'|N'}}(\Omega_{A'})$} pour $A'$,
en \unsure{$=$|$\simeq$} temps que $\Omega_{A'}$]. Car tout \unsure{hom}
% archive sheet number: 26

\scanpage{28}
\begingroup
\makeatletter\tagsleft@true\makeatother
à défaut d'appliquer le th. de dualité à $A$,
l'appliquer à $A'$, on trouve
\[
H^i(M)\simeq D\bigl(\operatorname{Ext}_{A'}^{\unsure{n'|N'}-i}(M,\Omega_{A'})\bigr)\qquad\text{i.e.}\tag*{(\deleted{\unsure{4}}1)}
\]
\[
\deleted{\operatorname{Ext}_{A'}^{\unsure{n'-i}}\unsure{k|(}}\tag*{\deleted{(1 bis)}}
\]
On en tire en particulier
\[
H^n(A)\simeq D\bigl(\operatorname{Ext}_{A'}^{n'-n}(M,\Omega_{A'})\bigr)\tag*{(2)}
\]
Ce qui prouve qu'on peut \deleted{\unsure{prendre}}\inserted{poser}
\[
\Omega_A\simeq\operatorname{Ext}_{A'}^{n'-n}(M,\Omega_{A'}).\tag*{(2 bis)}
\]

\deleted{Plus généralement, \unsure{introduisons}}

\[
\deleted{\Omega}\tag*{\deleted{(3)}}
\]
Soit $p$ un idéal premier de
$A$, \unsure{donc} $p\simeq A'/p'$, où $p'\supset\unsure{\mathfrak N|N}$. On a
alors
\[
(\Omega_A)_p\simeq\operatorname{Ext}_{A'_{p'}}^{n'-n}
(M_{p^{*}},(\Omega_{A'})_{p'}).
\]

Supposons qu'on sache que $(\Omega_{A'})_{p'}$ est
un module fond. pour l'anneau
[\uline{Coh. Mac.}] $A'_{p'}$ : ce sera le cas en
particulier si $A'$ est \uline{régulier} car alors
$\Omega_{A'}\simeq A'$, et $A'_{p'}$ sera régulier (Serre)
donc $A'_{p'}$ en sera un mod. fond. Alors
$(\Omega_A)_p$ est un mod. dualisant pour $\deleted{\unsure{A}}A_p$
d'après (2 bis) appliqué à $A_p\unsure{\Rightarrow|\simeq|=}A'_{p'}/\mathfrak N A'_{p'}$.
\endgroup
% archive sheet number: 27

\scanpage{29}
\begingroup
\makeatletter\tagsleft@true\makeatother
D'où

\uline{Prop 11} Soit $A$ un quotient d'un anneau
local régulier, \deleted{\unsure{si}} $\Omega$ est un module fond.
pour $A$. Si $p\subset A$ est premier \unsure{de|dans} $A$,
$\Omega_p$ est fond. pour $A_p$.

[N.B, Une démonstration directe à partir
des définitions semble \unsure{impratiquable|impraticable}]. De plus,
on tire de (1) le

\uline{Lemme} Pour que $\operatorname{\unsure{\deleted{co}codim|cocodim}}M\leq k$, il f. et
s. que $\operatorname{Ext}_{A'}^{\unsure{n'|N'}-i}(M,\Omega_{A'})=0$ pour $0\leq i<k$.

[Fait bien connu, dû à \unsure{Serre}]. \deleted{\unsure{\S}}

\uline{Corollaire} $M$ Coh-Mac $\Longleftrightarrow\operatorname{Ext}^{n'-n+i}(M,\Omega_{A'})=0$
\deleted{pour} $i>0$.

Cela nous amène à introduire les
\[
\Omega_A^{(i)}=\operatorname{Ext}_{A'}^{n'-n+i}(M,\Omega_{A'})\tag*{(3)}
\]
définis (à isom \uline{non can} près) par
\[
\deleted{\unsure{E}}\,\unsure{D|\mathbb D}(\Omega_A^{(i)})\simeq H^{\unsure{n-i}}(A)\tag*{(3 bis)}
\]
[Nous étudierons plus \unsure{loin} la \unsure{covariance|variance} exacte des $\Omega_A^{(i)}$.].
Donc $\Omega_A^{(0)}=\Omega_A$ mod. \unsure{fond|fond.}, et
\[
\Omega_A^{(i)}=0\quad\text{pour}\quad i\geq1\qquad\Longleftrightarrow\quad A\text{ est Coh. Mac.}
\]
\endgroup
% archive sheet number: 28

